\documentclass[12pt, a4paper]{report}

%---- GENERAL ---%
\usepackage[utf8x]{inputenc} % Encodage du fichier
\usepackage[T1]{fontenc}  % Encodage des polices du document
\usepackage[english, francais]{babel} % Langages utilisés
\usepackage{lmodern} % Police pour rendre le texte sélectionable

\usepackage[top=2.5cm, bottom=2.5cm, left=2.5cm, right=2.5cm]{geometry} % Marges PDF 

%---- TABLE DES MATIERES - PRESENTATION ---%
\usepackage[titles]{tocloft} % Package pour paramtérer la table des matières
\renewcommand{\cftchapleader}{\cftdotfill{\cftdotsep}} % Mets des points pour les chapitre et les sections
\usepackage[colorlinks=true,urlcolor=blue,pdfstartview=FitH,linkcolor=blue,hyperfootnotes=true]{hyperref} % Liens

%---- GLOSSAIRE ---%
\usepackage{glossaries}
\makeglossaries

%---- STRUCTURES DE DONNEES & MISE EN FORME ---%
\usepackage{enumerate} % Listes
\usepackage{graphicx} % Figures
\usepackage{color} % Couleur
\definecolor{darkblue}{RGB}{0,0,190}
\definecolor{lightgray}{gray}{0.70}
\definecolor{darkgreen}{RGB}{0,150,50}

%---- HEADERS & FOOTERS ---%
\usepackage{fancyhdr}
\pagestyle{fancy}
\renewcommand{\sectionmark}[1]{\markright{#1}}
\fancyhf{}

\fancyhead[RE]{\chaptername{} \thechapter{} - \nouppercase{\rightmark}} % Even page - Right side
\fancyhead[LE]{\thepage} % Even page - Left side
\fancyhead[RO]{\thepage}% Odd page - Rigth side
\fancyhead[LO]{\chaptername{} \thechapter{} - \nouppercase{\rightmark}} % Odd page - Left side
\fancyfoot[C]{footer} % Footer - Center

\newcommand{\clearemptydoublepage}{\newpage{\pagestyle{empty}\cleardoublepage}} % Supprime les en-têtes et pieds de page des pages vides (type book)

%---- CHAPTERS ---%
\usepackage[Lenny]{fncychap}

%---- CODES SOURCES & ALGORITHMES ---%
\usepackage{listings} % Codes sources

% [style = sourceC] for the C codes
\lstdefinestyle{sourceC}{language=C, basicstyle=\small\color{black}\ttfamily, classoffset=0, keywordstyle=\color{blue}\bfseries, classoffset=1, morekeywords={NULL}, keywordstyle=\color{magenta}\bfseries, commentstyle=\color{brightblue}\itshape, stringstyle=\color{darkgreen}, showstringspaces=false, tabsize=3, framexleftmargin=2mm, frame=shadowbox, rulesepcolor=\color{lightgray}, breaklines=true, emph={main, printf, scanf, FILE, fopen, fscanf, fprintf, fclose, rand}, emphstyle=\color{darkblue}\bfseries,columns=
fullflexible, flexiblecolumns=true, upquote=true, keepspaces=true}

% [style = inlineSourceC] for C expression into text
\lstdefinestyle{inlineSourceC}{language=C, basicstyle=\footnotesize\color{black}\ttfamily, classoffset=0, keywordstyle=\color{blue}\bfseries, classoffset=1, morekeywords={NULL}, keywordstyle=\color{magenta}\bfseries, commentstyle=\color{brightblue}\itshape, stringstyle=\color{darkgreen}, showstringspaces=false, tabsize=3, breaklines=true, emph={main, printf, scanf, FILE, fopen, fscanf, fprintf, fclose, rand}, emphstyle=\color{darkblue}\bfseries,columns=fullflexible, keepspaces=true,upquote=true}

% [style = sourceJava] for the Java codes
\lstdefinestyle{sourceJava}{language=Java, basicstyle=\small\color{black}\ttfamily, classoffset=0, keywordstyle=\color{darkgreen}\bfseries, commentstyle=\color{brightblue}\itshape, stringstyle=\color{darkgreen}, showstringspaces=false, tabsize=3, framexleftmargin=2mm, frame=shadowbox, rulesepcolor=\color{lightgray}, breaklines=true, emph={main, printf, FILE, fopen, fscanf, fprintf, fclose,new,return,this,super,null,break,continue,if,else,switch,case,default,do,while,for,break,new,return}, emphstyle=\color{darkred}\bfseries,columns=fullflexible,upquote=true}

% [style = sourceJava] for the Java codes
\lstdefinestyle{inlineSourceJava}{style=sourceJava, basicstyle=\normalsize\color{black}\ttfamily}

% [style = msgTerminal] for shell codes on black background
\lstdefinestyle{msgTerminal}{language=sh, basicstyle=\small\color{white}\ttfamily, keywordstyle=\color{white}, commentstyle=\color{white}\itshape, stringstyle=\color{white}, showstringspaces=false, framexleftmargin=3mm, xleftmargin=3mm, frame=none, tabsize=3, backgroundcolor=\color{black}, rulecolor=\color{black}, breaklines=true,columns=fullflexible,upquote=true}

% [style = msgTerminalW] for shell codes on white background
\lstdefinestyle{msgTerminalW}{language=sh, basicstyle=\small\color{black}\ttfamily, keywordstyle=\color{black}, commentstyle=\color{black}\itshape, stringstyle=\color{black}, showstringspaces=false, framexleftmargin=3mm, xleftmargin=3mm, frame=none, tabsize=3, backgroundcolor=\color{white}, rulecolor=\color{white}, breaklines=true,columns=fullflexible,upquote=true}

% [style = inlineTerminal] for shell expression into text
\lstdefinestyle{inlineTerminal}{language=sh, basicstyle=\footnotesize\color{black}\ttfamily\bfseries, keywordstyle=\color{black}, commentstyle=\color{black}\itshape, stringstyle=\color{black}\itshape, showstringspaces=false, tabsize=3, breaklines=true,upquote=true}


\usepackage[french]{algorithm2e} % Algorithmes 
\newenvironment{algo}[0]{\vrule~\vrule\itshape \begin{algorithm}[H] }{\normalfont\end{algorithm}}
\newcommand{\var}[1]{\textnormal{\texttt{#1}}}
\RestyleAlgo{plain}
\SetKwSwitch{Selon}{Cas}{Autre}{selon}{faire}{cas où}{autres cas}{fin}{fin} %Si la compilation bloque à algorithm2e, tenter la ligne du dessous à la place de celle-là
\SetKwSwitch{Selon}{Cas}{Autre}{selon}{faire}{cas où}{autres cas}{fin}0
%\SetKwRepeat{Repeter}{r\'ep\'eter}{tant que}
\SetKwFor{uTq}{tant que}{faire}{} % Tant que sans fin (équivaut a uSi)
\SetKwFor{oTq}{}{}{fin} % Tant que déjà ouvert, donc sans balise de depart
\SetKwIF{oSi}{}{}{}{}{}{}{fin}
\SetKwIF{oFn}{}{}{fonction}{}{}{}{fin}
\SetKwFunction{KwFn}{Fn}


%---- DOCUMENT ---%
% COMPILATION : 
% pdflatex squelette_rapport.tex
% bibtex squelette_rapport
% makeindex squelette_rapport
% pdflatex squelette_rapport.tex
% pdflatex squelette_rapport.tex
\title{\textbf{Projet d'exploration : détection des distances via une ESPCam} \\ 
\Large{--- Épita --- \\ 
\textit{Équipe NARA}} \\ 
\vspace{1,5cm} 
\includegraphics[width=0.3\textwidth]{./img/logo_epita} \includegraphics[width=0.2\textwidth]{./img/logo_equipe_robotique_exploration} 
\vspace{1,5cm}}
\date{\today}
\author{Nathan PODESTA 
\and Adrien MORVAN 
\and Romain PATURET 
\and Auguste HATTON}

\begin{document}
	\maketitle
	%-----------------------
        \chapter*{Résumé}
	    \addcontentsline{toc}{chapter}{Résumé}
	%-----------------------
Notre projet a pour but la détection de cibles sur une image capturée et le calcul de la distance entre la caméra et la cible détecté, qui peut être utilisé dans tout type d'environnement.
Pour ce faire, nous avons choisi d'utiliser une ESPCam qui va capter les images et les transférer via le réseau Wireless Fidelity, ensuite nous avons fait le choix d'utiliser le logiciel OpenCV afin de traiter le flux vidéo et de détecter la cible.
	
	%-----------------------
	% TABLE DES MATIERES %
	\tableofcontents
	%-----------------------

	%-----------------------%
	%     INTRODUCTION      %
	%-----------------------%
	
	%-----------------------
	\chapter{Introduction}
	%-----------------------

Docteur Loica AVANTHEY et Docteur Laurent BAUDOIN, chercheurs au sein du SEAL, le laboratoire de recherche de l'Epita, nous ont confié pour mission de construire un projet qui consiste à détecter une cible et à mesurer la distance qui la sépare de notre caméra. Nous avons choisi d'utiliser le logiciel de traitement d'image en temps réel OpenCV et le composant ESPCam.
\newline

Exemple tableau \& table : 

\begin{tabular}{|l|c|r|}
	\hline \textbf{case 00} & \textbf{case 01} & \textbf{case 03} \\
	\hline case 10 & case 11 & case 13 \\ 
	case 20 & case 21 & case 23 \\ 
	case 20 & case 21 & case 33 \\ 
	\hline
\end{tabular}

\begin{table}
\begin{center}
	\begin{tabular}{|c|c|}
		\hline case 00 & case 01 \\
		\hline case 10 & case 11 \\ 
		\hline 
	\end{tabular}
	\caption{Légende du tableau}
\end{center}
\end{table}


Sed ut perspiciatis unde omnis iste natus error sit voluptatem accusantium doloremque laudantium, totam rem aperiam, eaque ipsa quae ab illo inventore veritatis et quasi architecto beatae vitae dicta sunt explicabo. Nemo enim ipsam voluptatem quia voluptas sit aspernatur aut odit aut fugit, sed quia consequuntur magni dolores eos qui ratione voluptatem sequi nesciunt. Neque porro quisquam est, qui dolorem ipsum quia dolor sit amet, consectetur, adipisci velit, sed quia non numquam eius modi tempora incidunt ut labore et dolore magnam aliquam quaerat voluptatem. 

Test des listes : 

\begin{itemize}
	\item Item 1
	\item Item 2
	\item Item 3
	\item ...
\end{itemize}


Ut enim ad minima veniam, quis nostrum exercitationem ullam corporis suscipit laboriosam, nisi ut aliquid ex ea commodi consequatur? Quis autem vel eum iure reprehenderit qui in ea voluptate velit esse quam nihil molestiae consequatur, vel illum qui dolorem eum fugiat quo voluptas nulla pariatur?

\begin{itemize}
	\item[$\Rightarrow$] Item 1
	\item[$\bullet$] Item 2
	\item[$\circ$] Item 3
	\item[$\times$] Item 4
	\item[$\surd$] ...
\end{itemize}

At vero eos et accusamus et iusto odio dignissimos ducimus qui blanditiis praesentium voluptatum deleniti atque corrupti quos dolores et quas molestias excepturi sint occaecati cupiditate non provident, similique sunt in culpa qui officia deserunt mollitia animi, id est laborum et dolorum fuga. Et harum quidem rerum facilis est et expedita distinctio. 

\begin{enumerate}
	\item Item 1
	\item Item 2
	\item ...
\end{enumerate}

Nam libero tempore, cum soluta nobis est eligendi optio cumque nihil impedit quo minus id quod maxime placeat facere possimus, omnis voluptas assumenda est, omnis dolor repellendus. Temporibus autem quibusdam et aut officiis debitis aut rerum necessitatibus saepe eveniet ut et voluptates repudiandae sint et molestiae non recusandae. 

\begin{enumerate}[a)]
	\item Item 1
	\item Item 2
	\item ...
\end{enumerate}

Itaque earum rerum hic tenetur a sapiente delectus, ut aut reiciendis voluptatibus maiores alias consequatur aut perferendis doloribus asperiores repellat.



	%-----------------------%
	%    ETAT DE L'ART      %
	%-----------------------%

	%-----------------------
	\chapter{État de l'Art}
	\thispagestyle{empty}
	%-----------------------
	
Introduction de l'état de l'art - Lorem ipsum dolor sit amet, consectetur adipisicing elit, sed do eiusmod tempor incididunt ut labore et dolore magna aliqua. Ut enim ad minim veniam, quis nostrud exercitation ullamco laboris nisi ut aliquip ex ea commodo consequat. Duis aute irure dolor in reprehenderit in voluptate velit esse cillum dolore eu fugiat nulla pariatur. Excepteur sint occaecat cupidatat non proident, sunt in culpa qui officia deserunt mollit anim id est laborum.

Utilisation d'une référénce~\cite{NOM00}.


Ut enim ad minima veniam, quis nostrum exercitationem ullam corporis suscipit laboriosam, nisi ut aliquid ex ea commodi consequatur? Quis autem vel eum iure reprehenderit qui in ea voluptate velit esse quam nihil molestiae consequatur, vel illum qui dolorem eum fugiat quo voluptas nulla pariatur?

		%-----------------------
		\section{Premier axe}
		\label{premier_axe}
		%-----------------------

Utilisation d'une autre référénce~\cite{NOM05}. Et encore une autre~\cite{NOM12}. Ou bien plusieurs~\cite{WEB13, NOM09, NOM07}.
Lorem ipsum dolor sit amet, consectetur adipisicing elit, sed do eiusmod tempor incididunt ut labore et dolore magna aliqua. Ut enim ad minim veniam, quis nostrud exercitation ullamco laboris nisi ut aliquip ex ea commodo consequat. Duis aute irure dolor in reprehenderit in voluptate velit esse cillum dolore eu fugiat nulla pariatur. Excepteur sint occaecat cupidatat non proident, sunt in culpa qui officia deserunt mollit anim id est laborum.

Figure~:

\includegraphics[width=0.3\textwidth]{./img/logo_equipe_robotique_exploration}

Sed ut perspiciatis unde omnis iste natus error sit voluptatem accusantium doloremque laudantium, totam rem aperiam, eaque ipsa quae ab illo inventore veritatis et quasi architecto beatae vitae dicta sunt explicabo. Nemo enim ipsam voluptatem quia voluptas sit aspernatur aut odit aut fugit, sed quia consequuntur magni dolores eos qui ratione voluptatem sequi nesciunt. Neque porro quisquam est, qui dolorem ipsum quia dolor sit amet, consectetur, adipisci velit, sed quia non numquam eius modi tempora incidunt ut labore et dolore magnam aliquam quaerat voluptatem. 

\begin{figure}
\begin{center}
	\includegraphics[width=0.3\textwidth]{./img/logo_equipe_robotique_exploration}
	\caption{Légende de l'image}
	\label{fig:image1}
\end{center}
\end{figure}


Ut enim ad minima veniam, quis nostrum exercitationem ullam corporis suscipit laboriosam, nisi ut aliquid ex ea commodi consequatur? Quis autem vel eum iure reprehenderit qui in ea voluptate velit esse quam nihil molestiae consequatur, vel illum qui dolorem eum fugiat quo voluptas nulla pariatur?

At vero eos et accusamus et iusto odio dignissimos ducimus qui blanditiis praesentium voluptatum deleniti atque corrupti quos dolores et quas molestias excepturi sint occaecati cupiditate non provident, similique sunt in culpa qui officia deserunt mollitia animi, id est laborum et dolorum fuga. Et harum quidem rerum facilis est et expedita distinctio. Nam libero tempore, cum soluta nobis est eligendi optio cumque nihil impedit quo minus id quod maxime placeat facere possimus, omnis voluptas assumenda est, omnis dolor repellendus. Temporibus autem quibusdam et aut officiis debitis aut rerum necessitatibus saepe eveniet ut et voluptates repudiandae sint et molestiae non recusandae. Itaque earum rerum hic tenetur a sapiente delectus, ut aut reiciendis voluptatibus maiores alias consequatur aut perferendis doloribus asperiores repellat.
		
		%-----------------------
		\section{Second axe}
		%-----------------------

Référence interne, voir la section \ref{premier_axe} page~\pageref{premier_axe}. % fait référence à l'endroit où on a posé le label correspondant au nom (\label{premier_axe})
Lorem ipsum dolor sit amet, consectetur adipisicing elit, sed do eiusmod tempor incididunt ut labore et dolore magna aliqua. Ut enim ad minim veniam, quis nostrud exercitation ullamco laboris nisi ut aliquip ex ea commodo consequat. Duis aute irure dolor in reprehenderit in voluptate velit esse cillum dolore eu fugiat nulla pariatur. Excepteur sint occaecat cupidatat non proident, sunt in culpa qui officia deserunt mollit anim id est laborum.

\begin{lstlisting}[style=sourceC]
#include <stdio.h>

int main(void) {
	printf("Hello World!\n");
	return 0;
}
\end{lstlisting}		

 
Sed ut perspiciatis unde omnis iste natus error sit voluptatem accusantium doloremque laudantium, totam rem aperiam, eaque ipsa quae ab illo inventore veritatis et quasi architecto beatae vitae dicta sunt explicabo. Nemo enim ipsam voluptatem quia voluptas sit aspernatur aut odit aut fugit, sed quia consequuntur magni dolores eos qui ratione voluptatem sequi nesciunt. 

Code \lstinline[style=sourceC]!#include <stdio.h>! en ligne.


Ut enim ad minima veniam, quis nostrum exercitationem ullam corporis suscipit laboriosam, nisi ut aliquid ex ea commodi consequatur? Quis autem vel eum iure reprehenderit qui in ea voluptate velit esse quam nihil molestiae consequatur, vel illum qui dolorem eum fugiat quo voluptas nulla pariatur?

Code en provenance d'un fichier~:

\lstinputlisting[style=sourceC]{src/helloWorld.c} 

At vero eos et accusamus et iusto odio dignissimos ducimus qui blanditiis praesentium voluptatum deleniti atque corrupti quos dolores et quas molestias excepturi sint occaecati cupiditate non provident, similique sunt in culpa qui officia deserunt mollitia animi, id est laborum et dolorum fuga. Et harum quidem rerum facilis est et expedita distinctio. Nam libero tempore, cum soluta nobis est eligendi optio cumque nihil impedit quo minus id quod maxime placeat facere possimus, omnis voluptas assumenda est, omnis dolor repellendus. 

\begin{algorithm}[H]	
	Une instruction \\
	\Tq{La condition est vraie}{
		\eSi{Une autre condition est vraie}{
			Une instruction
		}{
			Une autre instruction
		}
	}
	\Si{Une condition est vraie}{
		Une instruction \\
		Une autre instruction \\	
	}
	\Repeter{La condition est vraie}{
		Une instruction
	}
	
\end{algorithm}


Temporibus autem quibusdam et aut officiis debitis aut rerum necessitatibus saepe eveniet ut et voluptates repudiandae sint et molestiae non recusandae. Itaque earum rerum hic tenetur a sapiente delectus, ut aut reiciendis voluptatibus maiores alias consequatur aut perferendis doloribus asperiores repellat.
			

Considérez par exemple l'équation \ref{eq:equation1} page~\pageref{eq:equation1}~:

\begin{equation}
\left(\frac{x^2}{y^3}\right)
\label{eq:equation1}
\end{equation}

\clearemptydoublepage


			
	%-----------------------%
	% ETUDE DE FAISABILITE  %
	%-----------------------%

	%-----------------------
	\chapter{Étude de faisabilité}
	%-----------------------

Introduction de l'étude de faisabilité - Lorem ipsum dolor sit amet, consectetur adipisicing elit, sed do eiusmod tempor incididunt ut labore et dolore magna aliqua. Ut enim ad minim veniam, quis nostrud exercitation ullamco laboris nisi ut aliquip ex ea commodo consequat. Duis aute irure dolor in reprehenderit in voluptate velit esse cillum dolore eu fugiat nulla pariatur. Excepteur sint occaecat cupidatat non proident, sunt in culpa qui officia deserunt mollit anim id est laborum.

	Sed ut perspiciatis unde omnis iste natus error sit voluptatem accusantium doloremque laudantium, totam rem aperiam, eaque ipsa quae ab illo inventore veritatis et quasi architecto beatae vitae dicta sunt explicabo. Nemo enim ipsam voluptatem quia voluptas sit aspernatur aut odit aut fugit, sed quia consequuntur magni dolores eos qui ratione voluptatem sequi nesciunt. Neque porro quisquam est, qui dolorem ipsum quia dolor sit amet, consectetur, adipisci velit, sed quia non numquam eius modi tempora incidunt ut labore et dolore magnam aliquam quaerat voluptatem.
	
		%-----------------------
		\section{Première catégorie}
		%-----------------------

Introduction de la première catégorie - Ut enim ad minima veniam, quis nostrum exercitationem ullam corporis suscipit laboriosam, nisi ut aliquid ex ea commodi consequatur? Quis autem vel eum iure reprehenderit qui in ea voluptate velit esse quam nihil molestiae consequatur, vel illum qui dolorem eum fugiat quo voluptas nulla pariatur?
		
			%-----------------------
			\subsection{Argumentation 1}
			%-----------------------

Lorem ipsum dolor sit amet, consectetur adipisicing elit, sed do eiusmod tempor incididunt ut labore et dolore magna aliqua. Ut enim ad minim veniam, quis nostrud exercitation ullamco laboris nisi ut aliquip ex ea commodo consequat. Duis aute irure dolor in reprehenderit in voluptate velit esse cillum dolore eu fugiat nulla pariatur. Excepteur sint occaecat cupidatat non proident, sunt in culpa qui officia deserunt mollit anim id est laborum.

Sed ut perspiciatis unde omnis iste natus error sit voluptatem accusantium doloremque laudantium, totam rem aperiam, eaque ipsa quae ab illo inventore veritatis et quasi architecto beatae vitae dicta sunt explicabo. Nemo enim ipsam voluptatem quia voluptas sit aspernatur aut odit aut fugit, sed quia consequuntur magni dolores eos qui ratione voluptatem sequi nesciunt. Neque porro quisquam est, qui dolorem ipsum quia dolor sit amet, consectetur, adipisci velit, sed quia non numquam eius modi tempora incidunt ut labore et dolore magnam aliquam quaerat voluptatem. 

Ut enim ad minima veniam, quis nostrum exercitationem ullam corporis suscipit laboriosam, nisi ut aliquid ex ea commodi consequatur? Quis autem vel eum iure reprehenderit qui in ea voluptate velit esse quam nihil molestiae consequatur, vel illum qui dolorem eum fugiat quo voluptas nulla pariatur?


			%-----------------------
			\subsection{Argumentation 2}
			%-----------------------


Sed ut perspiciatis unde omnis iste natus error sit voluptatem accusantium doloremque laudantium, totam rem aperiam, eaque ipsa quae ab illo inventore veritatis et quasi architecto beatae vitae dicta sunt explicabo. Nemo enim ipsam voluptatem quia voluptas sit aspernatur aut odit aut fugit, sed quia consequuntur magni dolores eos qui ratione voluptatem sequi nesciunt. Neque porro quisquam est, qui dolorem ipsum quia dolor sit amet, consectetur, adipisci velit, sed quia non numquam eius modi tempora incidunt ut labore et dolore magnam aliquam quaerat voluptatem. 

Ut enim ad minima veniam, quis nostrum exercitationem ullam corporis suscipit laboriosam, nisi ut aliquid ex ea commodi consequatur? Quis autem vel eum iure reprehenderit qui in ea voluptate velit esse quam nihil molestiae consequatur, vel illum qui dolorem eum fugiat quo voluptas nulla pariatur?




	%-----------------------%
	%       CONCLUSION      %
	%-----------------------%
	
	%-----------------------
	\chapter{Conclusion}
	%-----------------------


	
	%-----------------------%
	%      GLOSSAIRE        %
	%-----------------------%

\deftranslation{Glossary}{Glossaire}
\addcontentsline{toc}{chapter}{Glossaire} % Rajoute la partie dans le sommaire
\printglossaries
	
	%-----------------------%
	%      BIBLIOGRAPHIE    %
	%-----------------------%

\addcontentsline{toc}{chapter}{Bibliographie} % Rajoute la partie dans le sommaire
\bibliographystyle{alpha}
\bibliography{biblio/bibliographie.bib}

	
\end{document}